\pdfvariable trailerid {[<C3232026090300000000000000000000><C3232026090300000000000000000000>]}
\documentclass[10pt]{article}
\usepackage[margin=0.72in]{geometry}
\usepackage{amsmath,amssymb,amsthm,booktabs,microtype,hyperref}
\hypersetup{colorlinks=true,linkcolor=blue,urlcolor=blue}
\providecommand{\CRevisionRound}{2}
\newtheorem{theorem}{Theorem}
\newtheorem{corollary}[theorem]{Corollary}
\newcommand{\ket}[1]{|#1\rangle}
\newcommand{\bra}[1]{\langle#1|}
\newcommand{\iu}{\mathrm{i}}
\newcommand{\spec}{\operatorname{spec}}
\setlength{\emergencystretch}{3em}
\title{Exact Resonance and Detuning Geometry in Complete-Graph Quantum Search}
\author{Route-A source-local certificate HCS-C323}
\date{3 September 2026}

\begin{document}
\maketitle
\begin{abstract}
For continuous-time search with $M$ marked vertices among $N$, we retain the
entire Hilbert space rather than only the initialized two-level reduction.  A
bright--dark decomposition gives every eigenvalue and multiplicity and an
exact success law for arbitrary nonnegative driver strength.  In the
nontrivial marked-fraction regime, perfect search occurs if and only if the
driver and oracle are resonant.
\ifnum\CRevisionRound>0
The off-resonance defect yields a critical detuning window, while a scalar
shift closes the conversion to complete-graph adjacency.  All missing-sector
faces are resolved directly.
\fi
\ifnum\CRevisionRound>1
Exact, high-precision, replay, and hostile evidence audit the convention
boundary.  The Hamiltonian is a natural quantization, but it carries no
target arithmetic data and Route A is rejected.
\fi
\end{abstract}

\section{Exact success law and full spectrum}
Let $\mathcal H=\mathbb C^N$ have its computational basis, let $W$ be a
marked set of size $M$, and let $P_W$ project onto its span.  Write
\begin{equation}\label{eq:model}
 H_g=-g\ket{s}\bra{s}-P_W,\qquad
 \ket{s}=N^{-1/2}\sum_{j=1}^{N}\ket{j},\qquad g\geq0.
\end{equation}
The oracle coefficient is one.  Success means projection onto the \emph{full}
marked subspace, not the probability of one selected marked basis vector.

Assume first that $N\geq2$ and $0<M<N$, and put $a=M/N$.  The normalized
uniform marked and unmarked vectors obey
\begin{equation}\label{eq:brightstate}
 \ket{s}=\sqrt a\ket{w}+\sqrt{1-a}\ket{r}.
\end{equation}
Let $D_W$ comprise marked vectors whose coordinates sum to zero, and define
$D_R$ analogously on the unmarked coordinates.

\begin{theorem}[Bright--dark spectral theorem]\label{thm:main}
The orthogonal decomposition
\begin{equation}\label{eq:decomp}
 \mathcal H=D_W\oplus D_R\oplus\operatorname{span}\{\ket w,\ket r\}
\end{equation}
reduces $H_g$.  Its dark spectral data are
\begin{center}
\begin{tabular}{c@{\qquad}c@{\qquad}c}
sector& eigenvalue& multiplicity\\ \hline
$D_W$&$-1$&$M-1$\\
$D_R$&$0$&$N-M-1$.
\end{tabular}
\end{center}
On the bright sector the two eigenvalues are
\begin{equation}\label{eq:eigenvalues}
 \lambda_\pm=-\frac{g+1}{2}\pm\frac{\Omega}{2},\qquad
 \Omega^2=(g-1)^2+4ga.
\end{equation}
Starting from $\ket s$, the marked-subspace probability is
\begin{equation}\label{eq:success}
 p_W(t)=a+\frac{4ga(1-a)}{\Omega^2}
 \sin^2\!\left(\frac{\Omega t}{2}\right).
\end{equation}
For $0<a<1$, perfect success occurs if and only if $g=1$.  Its first hitting
time at resonance is
\begin{equation}\label{eq:firsthit}
 t_{\rm hit}=\frac{\pi}{2\sqrt a}=\frac{\pi}{2}\sqrt{\frac NM}.
\end{equation}
For $g>0$, the first peak is at $t_{\rm peak}=\pi/\Omega$, and
\begin{align}
 p_{\max}&=a+\frac{4ga(1-a)}{(g-1)^2+4ga},\label{eq:pmax}\\
 1-p_{\max}&=\frac{(1-a)(g-1)^2}{(g-1)^2+4ga}.\label{eq:defect}
\end{align}
For $g=0$, $p_W(t)=a$ at every time.
\end{theorem}

\begin{proof}
Every vector in $D_W$ is orthogonal to $\ket s$ and lies in the range of
$P_W$, so $H_g=-I$ there.  Every vector in $D_R$ is annihilated by both
summands in~\eqref{eq:model}.  The two dimensions and~\eqref{eq:decomp}
follow immediately.  In the ordered bright basis $(\ket w,\ket r)$,
\begin{equation}\label{eq:brightmatrix}
 H_B=-\begin{pmatrix}
 1+ga&g\sqrt{a(1-a)}\\
 g\sqrt{a(1-a)}&g(1-a)
 \end{pmatrix}.
\end{equation}
Its trace is $-(g+1)$ and its determinant is $g(1-a)$.  The discriminant is
$(g+1)^2-4g(1-a)=\Omega^2$, proving~\eqref{eq:eigenvalues} and completing the
full multiplicity ledger.

Set $B=H_B+(g+1)I/2$.  Cayley--Hamilton gives
$B^2=\Omega^2I/4$, hence
\begin{equation}\label{eq:exp}
 e^{-\iu tH_B}=e^{\iu t(g+1)/2}
 \left[\cos\!\left(\frac{\Omega t}{2}\right)I
 -\frac{2\iu}{\Omega}\sin\!\left(\frac{\Omega t}{2}\right)B\right].
\end{equation}
Moreover $\bra wB\ket s=-(g+1)\sqrt a/2$.  The initial state has no marked
dark component, so its full marked probability is the squared bright
amplitude.  Equation~\eqref{eq:exp} therefore gives
\[
 a\cos^2(\Omega t/2)+\frac{a(g+1)^2}{\Omega^2}
 \sin^2(\Omega t/2),
\]
which equals~\eqref{eq:success} because
$(g+1)^2-\Omega^2=4g(1-a)$.  At a peak, direct subtraction yields
\eqref{eq:defect}.  Its denominator is positive for $0<a<1$, and the defect
vanishes exactly when $g=1$.  Then $\Omega=2\sqrt a$, proving
\eqref{eq:firsthit}.  At $g=0$ the mixing coefficient vanishes, completing
the exceptional statement.
\end{proof}

\ifnum\CRevisionRound>0
\section{Critical detuning window and boundary atlas}
The factored defect is stronger than a statement that resonance suffices: it
also fixes the scale on which inaccurate driving remains useful.

\begin{corollary}[Sharp detuning window]\label{cor:window}
Let $k\in\mathbb N$, take $M=1$, $N=k^2$ (so $a=k^{-2}$), and set
$g=1+c/k$, where $c\in\mathbb R$ is fixed and $k>|c|$.
Then
\begin{equation}\label{eq:windowomega}
 \Omega^2=\frac{c^2+4+4c/k}{k^2},
\end{equation}
and, as $k\to\infty$,
\begin{equation}\label{eq:window}
 p_{\max}\longrightarrow\frac{4}{c^2+4},\qquad
 \sqrt a\,t_{\rm peak}\longrightarrow\frac{\pi}{\sqrt{c^2+4}}.
\end{equation}
Thus the nontrivial high-fidelity window has width $O(\sqrt a)$ around
$g=1$; fixed nonzero detuning instead drives $p_{\max}$ to zero.
\end{corollary}
\begin{proof}
Substituting $a=k^{-2}$ and $g=1+c/k$ gives~\eqref{eq:windowomega}.
Substitution into~\eqref{eq:pmax} and $t_{\rm peak}=\pi/\Omega$ proves both
limits.  For fixed $g\ne1$, the numerator of the oscillatory increment is
$O(a)$ while the denominator tends to $(g-1)^2$.
\end{proof}

The usual complete-graph formulation has no hidden normalization discrepancy.
Indeed, for the adjacency matrix $A(K_N)$,
\begin{equation}\label{eq:adjacency}
 A(K_N)=N\ket s\bra s-I.
\end{equation}
If $g=\gamma N$, then
\begin{equation}\label{eq:shift}
 -\gamma A(K_N)-P_W=H_g+\gamma I,
 \qquad
 e^{-\iu t(-\gamma A-P_W)}=e^{-\iu\gamma t}e^{-\iu tH_g}.
\end{equation}
Thus every probability agrees, although absolute eigenvalue ledgers differ by
$\gamma$.  Equation~\eqref{eq:shift} is a global-phase equivalence, not an
equality of Hamiltonians.

\begin{theorem}[Every missing-sector face]\label{thm:faces}
For $N\geq1$ the boundary data are as follows.
\begin{enumerate}
\item If $M=0$, the spectrum is $-g$ once and $0$ with multiplicity $N-1$;
success is identically zero.
\item If $M=N$, the spectrum is $-(g+1)$ once and $-1$ with multiplicity
$N-1$; success is identically one.
\item These formulas include $N=1$: the orthogonal multiplicity is zero, so
no nonexistent dark vector is introduced.
\item If $0<M<N$ and $g=0$, the spectrum is $-1$ with multiplicity $M$ and
$0$ with multiplicity $N-M$, while success is the constant $M/N$.
\end{enumerate}
\end{theorem}
\begin{proof}
When $M=0$, $P_W=0$ and the rank-one projector $\ket s\bra s$ supplies the
first spectrum.  When $M=N$, $P_W=I$ supplies the second.  A one-dimensional
space has no orthogonal complement.  Finally $g=0$ leaves the oracle
projector alone, so marked and unmarked amplitudes cannot mix.
\end{proof}
\fi

\ifnum\CRevisionRound>1
\section{Evidence collision and Route boundary}
The machine-readable evidence enumerates all $N=2,\ldots,32$, every
$M=1,\ldots,N-1$, and
\[
 g\in\{0,1/4,1/2,1,3/2,2,4\},
\]
giving 3,472 exact interior rows.  It adds 28 critical-window witnesses and
256 boundary rows, for 75,296 audited scalar leaves.  An independently
written checker reconstructs the schema, both hashes of the locked Route-A
YAML, every exact row, and selected full-matrix spectra and exponentials at
72 digits: 77,169 checks.  SymPy closes 17,885 exact identities; two isolated
producer runs replay byte for byte; 41 repaired-hash, parser, nonfinite,
duplicate-key, and semantic mutations must fail.  The finite grid is a
regression certificate.  Theorems~\ref{thm:main} and~\ref{thm:faces} and
Corollary~\ref{cor:window} have the stated all-parameter domains.

Farhi and Gutmann introduced continuous-time Hamiltonian search and its
square-root clock in the rank-one marked setting~\cite{farhigutmann}.  The
present arbitrary-$M$, arbitrary-detuning closure is a source-local
reconstruction; no literature-priority claim is made.  It is also disjoint
from the workspace's C143 discrete-time coined five-cycle walk, C171
stochastic Ehrenfest hypercube, C183 random-transposition Markov operator,
C223 Jaynes--Cummings blocks, and C318 one-dimensional SSH chain.  The
distinguishing object here is a permutation-symmetric oracle Hamiltonian and
its exact resonance window.

Under \texttt{NO\_BAD\_EULER\_OR\_ROOT\_NUMBER}, the strict Route-A tuple is
\[
 (\mathrm{A0\_FAIL},\mathrm{A1\_WEAK},\mathrm{A2\_FAIL},
  \mathrm{A3\_FAIL},\mathrm{A4\_NATURAL\_QUANTIZATION}).
\]
The dynamics has a genuine physical unitary clock and source-native Hermitian
quantization, but no rational-prime payload, arithmetic primitive-orbit
ledger, Euler product, target divisor, counting law, functional equation, or
target zero correspondence.  In particular $H_g$ is not claimed to be a
Hilbert--P\'olya operator.  Route A is rejected and Route B remains locked.

\begin{thebibliography}{9}
\bibitem{farhigutmann}
E.~Farhi and S.~Gutmann,
``An analog analogue of a digital quantum computation,''
\emph{Physical Review A} \textbf{57} (1998), 2403--2406.
\href{https://doi.org/10.1103/PhysRevA.57.2403}{doi:10.1103/PhysRevA.57.2403};
\href{https://arxiv.org/abs/quant-ph/9612026}{quant-ph/9612026}.
\end{thebibliography}
\fi

\end{document}
